\documentclass[11pt]{article}
\usepackage[margin=1.1in]{geometry}
\usepackage{amsmath, amssymb, amsthm}
\usepackage{booktabs}
\usepackage{graphicx}
\usepackage{hyperref}
\usepackage{microtype}

\newtheorem{proposition}{Proposition}
\newtheorem{corollary}{Corollary}
\newcommand{\E}{\mathbb{E}}
\newcommand{\dkl}{D_{\mathrm{KL}}}
\newcommand{\dent}{D_{\mathrm{ent}}}
\newcommand{\Preal}{P_{\mathrm{real}}}
\newcommand{\Qbs}{Q_{\mathrm{bs}}}

\title{Relative-Entropy Risk Discounting in Growth-Optimal Selection:\\
Weekly Vertical Credit Spreads on the S\&P~100}
\author{Peter~J.~Mercurio,~PhD}
\date{October 2026}

\begin{document}
\maketitle

\begin{abstract}
In this fourth paper of our series on entropic portfolio optimization we take
the framework developed in the first three papers into production application:
the systematic selection of short-dated vertical credit spreads on S\&P~100
underlyings. We introduce the \emph{intrinsic} GROUND ratio, a per-candidate
form of the Growth Rate Over UNiform Divergence ratio of our second paper,
which scores each candidate trade as a Kelly expected value under an estimated
outcome measure, discounted exponentially in a relative entropy. That entropy is
$\dent = \dkl(\Qbs \,\|\, U_3) = \ln 3 - H(\Qbs)$, the divergence of the
\emph{market's} three-state outcome measure from the discrete uniform, which is
exactly the risk measure of the second paper applied per candidate. Two facts
make it the right quantity to charge. First, on this instrument
family $\dent$ is an almost exact monotone transform of the market's
partial-zone probability, hence of the market's expected move measured in
strike widths (Spearman $-0.999$ on $25{,}228$ eligible candidates): the
entropy factor is a measurable statement about whether the spread is a genuine
three-state wager or a disguised binary one. Second, $\dent$ is a function of
quoted prices alone and is reproducible across independent data vendors at a
rank correlation of $0.98$, where the estimated win probability reproduces at
$0.54$; the penalty is therefore built from the robust half of the score and
the growth factor from the fragile half. We give the exact algebra of the
criterion, prove the monotonicity and second-order results, document the price
basis and the quote-reality gates that make a vendor backtest comparable to
live quotes, and report the configuration we call the fair credit model (FCM)
on $5{,}779$ in-sample
and $814$ out-of-time trades, together with a $146$-trade replay on live
brokerage quotes. We also report, and do not hide, the criterion's principal
measured limitation: GROUND orders the probability of \emph{losing} and does
not order dollars.
\end{abstract}

\noindent\textbf{Keywords:} intrinsic GROUND ratio; relative entropy; Shannon
entropy; Kullback--Leibler divergence; Kelly criterion; credit spreads;
variance risk premium; portfolio optimization; execution basis

\section{Introduction}
\label{sec:intro}

Systematic sellers of equity option premium face a selection problem rather
than an existence problem. That implied volatility exceeds subsequently
realized volatility on average is well established
\cite{carrwu2009, bollerslev2009, bakshikapadia2003}; that any particular
spread is cheap is not. On a single entry day in a hundred-name universe the
system described here scores on the order of a hundred candidate spreads, of
which roughly a fifth clear the liquidity and credit floors, and of those the
majority carry \emph{negative} model expected value. The task is to order the
remainder with a score whose ranking survives being computed on somebody
else's data.

This paper, the fourth in our series on portfolio optimization
\cite{mercurio2020repo, mercurio2020binary, mercurio2020gepo}, carries the
two-factor entropic structure of that series into a production configuration we
refer to throughout as the fair credit model, FCM. A growth factor prices each
candidate as a Kelly bet \cite{kelly1956} under an estimated physical measure,
an entropy factor measures a relative entropy in nats, and the product
discounts the first exponentially in the second. The substantive design question
is what that relative entropy should measure.

FCM takes it to be
\begin{equation}
\dent \;=\; \dkl(\Qbs \,\|\, U_3) \;=\; \ln 3 - H(\Qbs),
\label{eq:dent-intro}
\end{equation}
the divergence of the market's own three-state measure from the discrete
uniform $U_3$. This is not a new object: it is the risk measure of
\cite{mercurio2020binary}, where the relative entropy of a wager's outcome
distribution from the uniform reference is shown to be a convex risk measure
in the sense of F\"ollmer and Schied \cite{follmer2002}, with the uniform
distribution playing the role of the minimum-risk wager. Taking it per
candidate keeps the published form of the series and applies it to the measure
the market quotes rather than to the measure the trader estimates.

Three findings support the choice, and all three are reported with their
measurements in this paper.

\paragraph{The entropy factor is a statement about the instrument.} For a
vertical credit spread the market measure is
$\Qbs = (1 - \pi_s,\ \pi_\ell,\ \pi_s - \pi_\ell)$, where $\pi_s$ and
$\pi_\ell$ are risk-neutral exercise probabilities at the two strikes. The
short-leg delta band pins $\pi_s$ near one half, so $\dent$ varies almost
entirely with the partial-zone mass $\pi_s - \pi_\ell$, which is itself
determined by the market's expected move over the holding period measured in
strike widths. On the eligible pool the Spearman correlation between $\dent$
and the partial-zone mass is $-0.999$, and between $\dent$ and
$\sigma\sqrt{T}S/w_d$ is $+0.999$ (Section~\ref{sec:geometry}). A high-$\dent$
spread is one the market expects to blow through: the one-dollar ladder step
is narrow in its volatility units and the three-state wager has collapsed into
a binary one. Discounting such candidates is a statement about the instrument,
not about a forecast.

\paragraph{The entropy factor is the reproducible half of the score.} In a
reconciliation of the vendor end-of-day frame against independent brokerage
snapshots on the same spreads, $\dent$ reproduced at a rank correlation of
$0.98$ and the smile-fit model credit at $0.974$, while the estimated full-win
probability reproduced at only $0.54$. The score's fragility lives entirely in
the growth factor. Any penalty built on an estimated measure would carry that
fragility into the discount as well and compound it; a penalty taken from
quoted prices alone does not.

\paragraph{The two factors discriminate different things.} Measured as
within-day pairwise concordance on the eligible pool, $-\dent$ alone orders
the probability of \emph{not} taking a full loss at $0.560$ but orders dollars
at $0.480$, below chance; credit over width orders dollars at $0.582$ but
orders not-losing at $0.451$, below chance. Their product is the only score
tested that is above chance on both targets (Section~\ref{sec:concordance}).
The two-factor form is not decoration.

The paper also makes an empirical contribution independent of the choice of
entropy. A backtest of multi-leg option strategies is an
exercise in differencing quoted prices, and the quantities differenced are
small. Section~\ref{sec:basis} describes the price basis FCM uses -- a
per-chain smile fit, capped at the short leg's own risk-neutral exercise
probability and again at the best credit the quoted book ever showed -- and
the four quote-reality gates that remove candidates whose quoted books cannot
transact at all. Those gates cost the in-sample book $15\%$ of its profit and
$0.17$ of its Sharpe ratio, and we keep them, because a backtest that books
untransactable spreads is measuring its own data errors.

Section~\ref{sec:lit} reviews the related literature.
Section~\ref{sec:growth} develops the growth factor and the empirical belief
measure $\Preal$. Section~\ref{sec:entropy} develops $\dent$, its geometry and
its theory. Section~\ref{sec:criterion} assembles the criterion and states its
properties. Section~\ref{sec:basis} gives the price basis and the execution
model. Section~\ref{sec:design} gives the empirical design and
Section~\ref{sec:results} the results, including the ablations that isolate
each component's contribution, the fill ladder that locates the break-even
execution price, and the live-quote replay. Sections~\ref{sec:discussion}
and~\ref{sec:conclusions} discuss and conclude.

All logarithms are natural; entropies and divergences are in nats.

\section{Literature Review}
\label{sec:lit}

Portfolio selection begins with the mean-variance program of Markowitz
\cite{markowitz1952}, whose difficulties under non-Gaussian returns motivated
information-theoretic risk measures built on Shannon entropy
\cite{shannon1948}. Philippatos and Wilson \cite{philippatos1972} first
proposed entropy as a portfolio risk measure; Zhou, Cai, and Tong
\cite{zhou2013} survey the literature that followed.

The direct ancestors of the present system are the three papers of this
series. Return-entropy portfolio optimization (REPO)
\cite{mercurio2020repo} replaces the variance objective with portfolio entropy
computed combinatorially. Its extension to discrete-payoff assets
\cite{mercurio2020binary} is built on relative entropy and the Kelly
criterion: there the Kullback--Leibler divergence \cite{kullback1951} of a
wager's outcome distribution from the uniform reference is shown to be a
convex risk measure in the sense of F\"ollmer and Schied \cite{follmer2002},
with the uniform distribution itself the minimum-risk wager, and the
Growth Rate Over UNiform Divergence (GROUND) ratio is introduced as a
risk-adjusted comparison statistic in deliberate analogy to the
reward-to-variability ratio of Sharpe \cite{sharpe1966}. Generalized entropic
portfolio optimization (GEPO) \cite{mercurio2020gepo} extends the construction
to intervals of continuous returns and applies it to option strategies
including the vertical credit spreads traded here; its equation~19 is the
identity $\dkl(Q \| U_m) = \ln m - H(Q)$ that
Section~\ref{sec:entropy} uses directly. The present paper carries that
construction into production as a per-candidate score.

On the growth side, Kelly \cite{kelly1956} established the log-optimal
criterion, Thorp \cite{thorp1971} carried it to securities markets, Cover and
Thomas \cite{coverthomas2006} give the information-theoretic treatment in
which the growth forgone by betting a wrong distribution is a relative
entropy, and MacLean, Thorp, and Ziemba \cite{maclean2011} collect the
theory and practice including the fractional-Kelly variants used in
Section~\ref{sec:sizing}. On the robust-preferences side, the functional
$g - k\,D$ is the multiplier form of Hansen and Sargent
\cite{hansensargent2008}, a member of the variational class of Maccheroni,
Marinacci, and Rustichini \cite{maccheroni2006}, and the
exponential-in-divergence representation is the signature of entropic risk
measures \cite{follmer2002}. We are careful in
Section~\ref{sec:criterion} to state exactly how much of that identification
survives the change of reference measure, and how much does not.

Option-pricing inputs are standard: Black and Scholes
\cite{blackscholes1973} and Merton \cite{merton1973} for the $N(d_2)$
exercise probabilities, Breeden and Litzenberger \cite{breeden1978} for the
reading of those probabilities as state prices, Frittelli \cite{frittelli2000}
for the minimal-entropy martingale measure that is the incomplete-markets
limit of this family of criteria. On execution, Muravyev and Pearson
\cite{muravyev2020} show that effective option trading costs are far below
quoted spreads, the empirical fact that makes a midpoint-referenced basis
defensible; Christensen and Prabhala \cite{christensen1998} and Andersen
et al.\ \cite{andersen2003} characterize the implied--realized relation and the
noise in realized-volatility estimation that
Section~\ref{sec:discussion} returns to.

\section{The Growth Factor}
\label{sec:growth}

\subsection{The Kelly criterion}

Consider a repeated wager paying $+b$ per unit staked with probability $p$ and
$-1$ with probability $q = 1-p$. A bettor staking a constant wealth fraction
$w$ compounds at the expected exponential growth rate
\begin{equation}
G(w) \;=\; p \ln(1 + wb) + q \ln(1 - w),
\label{eq:kelly2}
\end{equation}
and the Kelly criterion selects the maximizer
$w^\ast = (pb - q)/b$, the edge divided by the odds \cite{kelly1956}.

\subsection{The three-state credit spread}
\label{sec:instrument}

A bull put credit spread sells a put at strike $K_s$ and buys a put at a lower
strike $K_\ell$; a bear call spread sells a call at $K_s$ and buys a call at a
higher strike. Write $c$ for the net credit per share, $w_d = |K_s - K_\ell|$
for the strike width, $m = w_d - c$ for the maximum loss per share, and
\begin{equation}
b \;=\; c / m
\label{eq:odds}
\end{equation}
for the payoff odds. At expiry the position resolves into one of three states:
a full win, in which the underlying settles on the safe side of $K_s$ and the
seller keeps $c$; a full loss, in which it settles beyond $K_\ell$ and the
seller pays $m$; and a partial outcome between the strikes. We write the
outcome measure as the triple $(p, q, r_0)$, with $p$ the full-win
probability, $q$ the full-loss probability, and $r_0 = 1 - p - q$ the
partial-zone mass.

With the per-unit-risk payoff vector $\{+b,\ +\alpha b,\ -1\}$, the
partial-zone coefficient follows from the linear payoff between the strikes
under a uniform location assumption:
\begin{equation}
\alpha(b) \;=\;
\begin{cases}
\dfrac{b-1}{2b}, & b < 1,\\[6pt]
0, & b \geq 1,
\end{cases}
\label{eq:alpha}
\end{equation}
so for $b < 1$ the partial state pays $\alpha b = (b-1)/2$ per unit risked,
the mean of $+b$ at the short strike and $-1$ at the long strike, which is
negative. For $b \geq 1$ setting the coefficient to zero understates the value
of partial outcomes and is retained as a conservative simplification.

The three-state growth rate is
\begin{equation}
\ell(w) \;=\; p \ln(1 + wb) + r_0 \ln(1 + w\alpha b) + q \ln(1 - w),
\label{eq:kelly3}
\end{equation}
whose first-order condition is a quadratic in $w$, linear when $\alpha = 0$
with root $w^\ast = (pb - q)/(b(p+q))$. We write $w^\ast$ for the maximizer,
clipped to $[0.01, 0.99]$, and define the Kelly expected value and the growth
signal as
\begin{equation}
E \;=\; e^{\ell(w^\ast)} - 1, \qquad g \;=\; \ln E .
\label{eq:kellyev}
\end{equation}
Because $\ell$ is the expectation of a concave function of the outcome, a
variance penalty is built into $E$ by Jensen's inequality: of two candidates
with equal arithmetic expected value, the higher-variance one has the lower
$E$.

\paragraph{Growth-negative candidates are evaluated, not discarded.} When the
first-order condition has no interior root the candidate is growth-negative
under the trader's measure. Such candidates are evaluated at the floor stake
$w = 0.01$ rather than returned as undefined. This matters more than it
appears: in an earlier revision of the system these rows were silently dropped
by the backtest while the live ranker retained them, which removed between
$38\%$ and $59\%$ of bull candidates per year from the backtest universe and
made the two environments incomparable. Under FCM, $60.6\%$ of the eligible
pool is growth-negative at the model credit and $15.9\%$ of selected picks
are; a criterion with no growth threshold must therefore be able to score
them.

\subsection{\texorpdfstring{$\Preal$}{P-real}: the empirical belief measure}
\label{sec:preal}

The trader's measure is not a parametric model. For each candidate, $\Preal$
is counted directly from the underlying's own realized history against the
candidate's own strikes.

Let $d = \min(\max(\mathrm{DTE}, 1), 4)$ be the holding period in sessions and
let $\{S_t\}$ be the name's daily closes. Form the $d$-session arithmetic
returns $R_t = S_{t+d}/S_t - 1$ over the $W = 252$ most recent windows that
\emph{completed strictly before} the entry date. Remove the window mean,
\begin{equation}
\tilde R_t \;=\; R_t - \frac{1}{W}\sum_{u} R_u ,
\label{eq:demean}
\end{equation}
add a per-candidate drift $\mu$ (Section~\ref{sec:drift}), and count crossings
of the two strike thresholds $\theta_s = K_s/S - 1$ and
$\theta_\ell = K_\ell/S - 1$, where $S$ is the entry spot. For a bull put,
writing $n_s = \#\{\tilde R_t + \mu \leq \theta_s\}$ and
$n_\ell = \#\{\tilde R_t + \mu \leq \theta_\ell\}$,
\begin{equation}
\Preal \;=\; \frac{\bigl(W - n_s + \gamma,\ \ n_\ell + \gamma,\ \
n_s - n_\ell + \gamma\bigr)}{W + 3\gamma},
\qquad \gamma = \tfrac12 ,
\label{eq:preal}
\end{equation}
with the inequalities reversed for a bear call. The pseudo-count
$\gamma = 1/2$ is the only regularization.

Three properties of (\ref{eq:preal}) deserve emphasis. It is strictly causal:
only completed windows enter. It is strike-exact: the thresholds are the
candidate's own strikes, not a delta proxy, so no interpolation between
instruments occurs. And it is \emph{drift-free} by construction
(\ref{eq:demean}): without the demeaning, the trailing year's realized drift
is baked into $p$, and the selection tilts mechanically toward the previous
year's winners. Under the raw-return version the median trailing one-year
return of selected names was $+31\%$ against $+3\%$ after demeaning. The cost
of removing it is that the surviving per-candidate edges are roughly an order
of magnitude smaller, which is why FCM operates with no growth threshold at
all (Section~\ref{sec:selection}).

\subsection{The own-gap drift}
\label{sec:drift}

The one quantity FCM permits to tilt $\Preal$ directionally is the
candidate's own opening gap on the entry day. Let
$\mathrm{gap} = (O_t / S_{t-1} - 1)$ divided by the prior 14-session Wilder
average true range in price units, let $z$ be its standardization clipped to
$\pm 4$ using a mean and standard deviation fit only on stock-days strictly
before the entry year, let $\sigma_d$ be the standard deviation of the name's
last 20 daily log close changes, and let $\beta_Y$ be a slope fit on the same
pre-$Y$ sample. Then
\begin{equation}
\mu \;=\; \beta_Y \cdot z \cdot \sigma_d \sqrt{d} .
\label{eq:drift}
\end{equation}
The fit is walk-forward and refit each January. No market average and no
cross-sectional factor enters: each name's drift comes from its own gap only.
A name with no gap recorded that day scores with $\mu = 0$. The fitted slopes
are small, from $0.113$ for entry year 2021 down to $0.012$--$0.026$
thereafter, so (\ref{eq:drift}) is a mild perturbation of $\Preal$ and not a
directional forecast in its own right.

\subsection{Kelly growth with carry}
\label{sec:carry}

Capital at risk in a credit spread earns financing over the holding period. Let
$\kappa = r \cdot \max(d, 1)/365$ with $r = 4\%$. The carry-adjusted growth
rate replaces the three payoffs of (\ref{eq:kelly3}) by $b + \kappa$,
$\alpha b + \kappa$ and $-(1 - \kappa)$:
\begin{equation}
\ell_\kappa(w) = p \ln\!\bigl(1 + w(b + \kappa)\bigr)
+ r_0 \ln\!\bigl(1 + w(\alpha b + \kappa)\bigr)
+ q \ln\!\bigl(1 - w(1 - \kappa)\bigr).
\label{eq:kellycarry}
\end{equation}
The shifted loss leg destroys the closed form, so $w^\ast_\kappa$ is taken on
a $400$-point grid. Two facts about (\ref{eq:kellycarry}) are worth recording.
It barely moves the stake: $\kappa \approx 2\times 10^{-4}$ against payoff
odds near unity, and books sized on $w^\ast_\kappa$ land within $1\%$ of books
sized on $w^\ast$. And it is exactly rank-neutral in selection: every entry
day in this design carries a single DTE, because Monday through Thursday
entries all expire the same Friday, so $\kappa$ is a per-day constant and
changes no within-day ordering. Selection therefore reads the uncarried score;
the Kelly sizing arms of Section~\ref{sec:sizing} read $w^\ast_\kappa$.

\section{The Entropy Factor}
\label{sec:entropy}

\subsection{Shannon entropy and relative entropy}

The Shannon entropy \cite{shannon1948} of a discrete distribution $P$ on
support $\mathcal{X}$ and the Kullback--Leibler divergence
\cite{kullback1951} of $P$ from $Q$ are
\begin{equation}
H(P) = -\sum_{x} P(x) \ln P(x), \qquad
\dkl(P \| Q) = \sum_{x} P(x) \ln \frac{P(x)}{Q(x)} .
\label{eq:hdkl}
\end{equation}
Gibbs' inequality (see MacKay \cite{mackay2003}) gives
$\dkl(P \| Q) \geq 0$ with equality if and only if $P = Q$ almost everywhere.

\subsection{The GROUND ratio and its intrinsic form}

To compare growth rates of portfolios with different risk profiles,
\cite{mercurio2020binary} introduced the Growth Rate Over UNiform Divergence
ratio. For a chosen portfolio $R_a$ and the minimum-risk portfolio $R_b$ in an
$m$-state event space with discrete uniform $U_m$,
\begin{equation}
G_m \;=\;
\frac{\E\bigl(G_a(\mathbf{w}_a) - G_b(\mathbf{w}_b)\bigr)}
     {\dkl(R_a \| U_m) - \dkl(R_b \| U_m)}
\;=\;
\frac{\E\bigl(G_a(\mathbf{w}_a) - G_b(\mathbf{w}_b)\bigr)}
     {H(R_b) - H(R_a)} .
\label{eq:ground-orig}
\end{equation}
The second equality is the identity
$\dkl(Q \| U_m) = \ln m - H(Q)$, equation~19 of \cite{mercurio2020gepo}, and
it is the form we use below.

Equation (\ref{eq:ground-orig}) is comparative: the score of a candidate
depends on a reference member of the same pool. In a streaming daily-selection
setting, where a handful of positions is rebuilt from scratch each day from a
changing candidate pool, that coupling is awkward. The \emph{intrinsic} GROUND
ratio replaces it with a per-candidate score requiring no reference member.
For a candidate with Kelly expected value $E$ from (\ref{eq:kellyev}), a
relative entropy $D$, and discount strength $k > 0$,
\begin{equation}
\Gamma \;=\; E \, e^{-kD} \;=\; \exp(g - kD) .
\label{eq:intrinsic}
\end{equation}

\subsection{\texorpdfstring{$\dent$}{D-ent}: the market's measure against the
uniform}
\label{sec:dent}

The form (\ref{eq:intrinsic}) is general in $D$. FCM takes
\begin{equation}
\boxed{\ \dent \;=\; \dkl(\Qbs \,\|\, U_3) \;=\; \ln 3 - H(\Qbs)\ }
\label{eq:dent}
\end{equation}
where $\Qbs$ is the market's three-state outcome measure for the same spread.
It is constructed from the quoted surface alone. For each chain
$(\text{name}, \text{date}, \text{expiry})$ a robust weighted quadratic is fit
to quoted implied volatilities in standardized moneyness
$y = \ln(K/S)/(\sigma_0\sqrt{T})$, with $\sigma_0$ the median implied
volatility of the six nearest strikes, weights
$1/(1 + \text{relative quoted width})$, puts and calls pooled, and three
rounds of $2.5$-MAD trimming. Writing $\hat\sigma_s$ and $\hat\sigma_\ell$ for
the fitted volatilities at the two strikes and
\begin{equation}
\pi_j \;=\; N\!\left(\frac{\ln(S/K_j) - \tfrac12\hat\sigma_j^2 T}
{\hat\sigma_j\sqrt{T}}\right)
\quad\text{(calls)}, \qquad
\pi_j \;=\; 1 - N(\cdot) \quad\text{(puts)},
\label{eq:nd2}
\end{equation}
the risk-neutral exercise probabilities \cite{blackscholes1973, merton1973},
with $\pi_\ell$ floored below $\pi_s$, the market triple is
\begin{equation}
\Qbs \;=\; \bigl(1 - \pi_s,\ \ \pi_\ell,\ \ \pi_s - \pi_\ell\bigr),
\label{eq:qbs}
\end{equation}
normalized to sum to one. By Breeden and Litzenberger \cite{breeden1978} these
are the state prices of the three outcomes up to discounting, so
(\ref{eq:qbs}) is the market's own probability statement about the position.

Two remarks on why the reference is $U_3$ and the argument is $\Qbs$.

The reference is $U_3$ because that is the published choice
\cite{mercurio2020binary, mercurio2020gepo}: $\dkl(\cdot \| U_m)$ is a convex
risk measure in the sense of \cite{follmer2002} on a finite outcome space, the
uniform distribution is its minimizer, and conventional risk measures --
variance, value at risk, conditional value at risk -- are defined on
continuous return distributions and become degenerate on a three-point
support. Relative entropy to the uniform fills exactly that gap, and does so
with the three-state wager's own geometry rather than a borrowed one.

The argument is $\Qbs$ for a reason this paper can measure rather than assert.
$\dent$ is a deterministic function of quoted prices, so it introduces no
estimation error of its own into the discount; the score's only estimated input
is the growth factor, and it is isolated there. The reconciliation study of
Section~\ref{sec:reproducibility} measures the consequence directly. On spreads matched by name, date and both strikes
across two independent data sources, $\dent$ reproduces at a rank correlation
of $0.98$ and the model credit at $0.974$, while the estimated full-win
probability reproduces at $0.54$. The score is a product of a fragile factor
and a robust one, and (\ref{eq:dent}) is the robust one.

\subsection{What \texorpdfstring{$\dent$}{D-ent} measures on this instrument}
\label{sec:geometry}

Because the short-leg delta band is narrow, $\pi_s$ is nearly constant across
the candidate pool and $\dent$ is driven by one quantity.

\begin{proposition}[Monotonicity in the partial-zone mass]
\label{prop:monotone}
Fix $\pi_s$ and write $m = \pi_s - \pi_\ell$ for the partial-zone mass, so
$\Qbs = (1 - \pi_s,\ \pi_s - m,\ m)$ with $0 < m < \pi_s$. Then $H(\Qbs)$ is
strictly increasing in $m$ for $m < \pi_s/2$ and strictly decreasing for
$m > \pi_s/2$; consequently $\dent$ is strictly \emph{decreasing} in $m$ on
$(0, \pi_s/2)$ and attains its minimum over $m$ at $m = \pi_s/2$.
\end{proposition}

\begin{proof}
Only the second and third components of $\Qbs$ depend on $m$. Write
$h(m) = -(\pi_s - m)\ln(\pi_s - m) - m \ln m$. Then
$h'(m) = \ln(\pi_s - m) - \ln m = \ln\bigl((\pi_s - m)/m\bigr)$, which is
positive for $m < \pi_s/2$, zero at $m = \pi_s/2$, and negative above. Since
$H(\Qbs) = -(1-\pi_s)\ln(1-\pi_s) + h(m)$ and $\dent = \ln 3 - H(\Qbs)$, the
claim follows.
\end{proof}

With $\pi_s \in [0.50, 0.60]$ the turning point $\pi_s/2$ lies in
$[0.25, 0.30]$, and on the FCM eligible pool only $1.8\%$ of candidates have
$m \geq 0.25$. Over the remaining $98.2\%$, $\dent$ is strictly decreasing in
the market's partial-zone mass. The measured relationship is almost perfectly
monotone: Spearman$(\dent, m) = -0.999$ on $25{,}228$ eligible candidates.

The partial-zone mass, in turn, is a statement about the market's expected
move relative to the strike ladder. Measuring that move in dollars as
$\hat\sigma_s \sqrt{T} S$ and dividing by the strike width gives
Spearman$(\dent,\ \hat\sigma_s\sqrt{T}S/w_d) = +0.999$ on the same pool. The
economic reading of (\ref{eq:dent}) follows:

\begin{quote}
$\dent$ is small when the market's expected move over the holding period is
small relative to the strike width, so the spread has a genuine third state
with real probability mass; $\dent$ is large when the expected move dwarfs the
width, so the spread has degenerated into a binary bet on the short strike.
The factor $e^{-k\dent}$ rewards the former and suppresses the latter.
\end{quote}

Table~\ref{tab:dentquint} shows the measured structure. Across $\dent$
quintiles the market's full-win probability barely moves, from $0.416$ to
$0.437$; what moves is the reallocation of mass from the partial state
($0.190 \to 0.055$) into the full-loss state ($0.388 \to 0.512$). High $\dent$
is not the market being bullish or bearish. It is the market pricing away the
cushion between the strikes.

\begin{table}[h]
\centering
\caption{The market measure $\Qbs$ by $\dent$ quintile on the FCM eligible
pool ($n = 25{,}228$); medians.}
\label{tab:dentquint}
\begin{tabular}{lrrrr}
\toprule
$\dent$ quintile & $\dent$ & $\Qbs$ full win & $\Qbs$ full loss & $\Qbs$ partial \\
\midrule
1 (lowest) & 0.0532 & 0.4157 & 0.3881 & 0.1897 \\
2          & 0.0950 & 0.4261 & 0.4294 & 0.1434 \\
3          & 0.1309 & 0.4322 & 0.4540 & 0.1144 \\
4          & 0.1688 & 0.4345 & 0.4775 & 0.0892 \\
5 (highest)& 0.2334 & 0.4366 & 0.5120 & 0.0553 \\
\bottomrule
\end{tabular}
\end{table}

A corollary is that $\dent$ carries a term-structure tilt. Longer holding
periods mean larger expected moves against the same ladder, so the median
$\dent$ rises with DTE: $0.1033$ at one day, $0.1190$ at two, $0.1327$ at
three, $0.1428$ at four. At a fixed $k$ the criterion therefore prefers short
holding periods, which it should: the same credit over fewer days is a higher
rate.

\subsection{Range and second-order form}

\begin{proposition}[Range and the quadratic approximation]
\label{prop:range}
On a three-state support, $0 \leq \dent \leq \ln 3$, with $\dent = 0$ if and
only if $\Qbs = U_3$. Moreover, as $\Qbs \to U_3$,
\begin{equation}
\dent \;=\; \tfrac32 \sum_{j=1}^{3}\Bigl(\Qbs(j) - \tfrac13\Bigr)^2
+ O(\delta^3), \qquad \delta = \max_j \bigl|\Qbs(j) - \tfrac13\bigr| ,
\label{eq:chi2}
\end{equation}
so to second order the penalty is $\tfrac32$ times the squared Euclidean
distance of the market's measure from equiprobability.
\end{proposition}

\begin{proof}
Non-negativity and the equality case are Gibbs' inequality; the upper bound is
attained in the limit as $\Qbs$ concentrates on one state, where
$H \to 0$. For (\ref{eq:chi2}) write $\Qbs(j) = \tfrac13 + \varepsilon_j$ with
$\sum_j \varepsilon_j = 0$. Then
$\dkl(\Qbs \| U_3) = \sum_j (\tfrac13 + \varepsilon_j)
\ln(1 + 3\varepsilon_j)
= \sum_j (\tfrac13 + \varepsilon_j)\bigl(3\varepsilon_j -
\tfrac92 \varepsilon_j^2\bigr) + O(\delta^3)
= \tfrac32 \sum_j \varepsilon_j^2 + O(\delta^3)$,
using $\sum_j \varepsilon_j = 0$. The leading term is
$\tfrac12 \chi^2(\Qbs, U_3)$ with
$\chi^2(\Qbs, U_3) = \sum_j (\Qbs(j) - \tfrac13)^2/\tfrac13$.
\end{proof}

The bounds in Proposition~\ref{prop:range} set the scale of $k$. The discount
factor lies in $[3^{-k}, 1]$, so at the canonical $k = 24$ it spans eleven
orders of magnitude in principle. In practice the realized range is far
narrower but still wide: on the eligible pool $\dent$ runs from $0.0069$ to
$0.3684$ with median $0.1309$, and the discount $e^{-24\dent}$ runs from
$0.0019$ at the fifth percentile to $0.409$ at the ninety-fifth, median
$0.043$. The entropy factor is not a tiebreaker; at $k = 24$ it is the
dominant term in the ordering, and Section~\ref{sec:ksweep} shows what
happens as $k$ varies. Note also that within the $[0.50, 0.60]$ delta band
$\dent$ cannot reach zero: the best attainable value, at $\pi_s = 0.60$ and
$m = 0.30$, is $\dent = 0.0097$, close to the observed minimum $0.0069$.

\section{The FCM Criterion}
\label{sec:criterion}

\subsection{Definition}

Collecting (\ref{eq:kellyev}), (\ref{eq:intrinsic}) and (\ref{eq:dent}), the
FCM score of a candidate is
\begin{equation}
\Gamma \;=\; \bigl(e^{\ell(w^\ast)} - 1\bigr)\, e^{-k\,\dent}
\;=\; E \, e^{-k(\ln 3 - H(\Qbs))}
\;=\; \exp\bigl(g - k\,\dent\bigr),
\qquad k = 24 .
\label{eq:fcm}
\end{equation}
The first factor is the variance-adjusted per-trade geometric expected return
at Kelly-optimal sizing under $\Preal$; the second is an entropic discount in
$(0, 1]$ computed from quoted prices. The product is dimensionless and reads
as an entropy-discounted per-trade geometric expected return. FCM applies
\emph{no threshold} to $\Gamma$:
selection is a credit floor followed by a daily cap on $\Gamma$-rank
(Section~\ref{sec:selection}). Section~\ref{sec:nothreshold} explains why,
and quantifies the consequence.

\subsection{Properties}

\begin{proposition}[Exponential form and the geometry of a $\Gamma$-cut]
\label{prop:halfplane}
Write $J_k = g - k\,\dent$. Then (i) $\Gamma = \exp(J_k)$ exactly, so for any
$\theta > 0$ the rule $\Gamma \geq \theta$ is identical to
$J_k \geq \ln\theta$ and the two induce the same ranking; and (ii) in the
$(\dent, g)$ plane the acceptance region of such a rule is the closed
half-plane
\begin{equation}
g \;\geq\; \ln\theta + k\,\dent ,
\label{eq:halfplane}
\end{equation}
a line of slope $k$ and intercept $\ln\theta$.
\end{proposition}

\begin{proof}
(i) $\exp(J_k) = e^g e^{-k\dent} = E e^{-k\dent} = \Gamma$ by
$g = \ln E$; rule equivalence follows from strict monotonicity of the
exponential. (ii) Take logarithms of $\Gamma \geq \theta$ and rearrange.
\end{proof}

The boundary of (\ref{eq:halfplane}) is the level set of the Lagrangian of
entropy-budgeted growth selection -- maximize $g$ subject to
$\dent \leq \varepsilon$ -- so $k$ has a direct reading as the shadow price of
market informational concentration, in nats of log-growth charged per nat of
the market measure's divergence from equiprobability. A daily top-$n$ rule is
the same half-plane with $\theta$ chosen each day so that exactly $n$
candidates clear it.

\begin{proposition}[Equivalent additive form]
\label{prop:additive}
$\ln\Gamma = g + k\,H(\Qbs) - k \ln 3$. Since $k\ln 3$ is a constant, FCM
ranking is identical to ranking on $g + k\,H(\Qbs)$: a linear trade-off
between the trader's log Kelly expected value and the market's Shannon
entropy, at $k$ nats of growth per nat of entropy.
\end{proposition}

\begin{proof}
Immediate from (\ref{eq:fcm}) and (\ref{eq:dent}).
\end{proof}

Proposition~\ref{prop:additive} is the cleanest statement of what FCM is. It
is a growth criterion that \emph{pays} for market uncertainty. This is the
reading the strategy was originally specified by: bet the direction the stock
is predicted to drift when that prediction beats the market's odds, favour
names where the market's view is high-entropy, and take only bets whose fill
yields positive expected value. The first clause is $g$ under $\Preal$ against
the market's credit, the second is $+k H(\Qbs)$, and the third is the
execution gate of Section~\ref{sec:execution}.

\subsection{What kind of preference \texorpdfstring{$\Gamma$}{Gamma} represents}
\label{sec:hs}

The functional $J_k = g - k\,\dent$ is linear in a relative entropy and so sits
in the variational class of Maccheroni, Marinacci, and Rustichini
\cite{maccheroni2006}, and the exponential-in-divergence form of
(\ref{eq:fcm}) is the signature of entropic risk measures
\cite{follmer2002}. We are precise about which member it is, because the
multiplier preferences of Hansen and Sargent \cite{hansensargent2008} have the
same shape and a different meaning.

The multiplier functional charges an agent for the distance of its \emph{own}
measure from a reference, so the penalty moves when the agent's beliefs move.
In (\ref{eq:fcm}) the divergence does not involve $\Preal$ at all. It is a
property of the quoted surface, identical for every agent looking at the same
chain, and an agent whose beliefs change does not change its own penalty.
$\Gamma$ is therefore not a robust-control criterion: it is a growth criterion
carrying an \emph{instrument-level} entropic risk charge, in the sense of
\cite{follmer2002} and of \cite{mercurio2020binary}, whose ambiguity index is
a state of the market rather than a property of the decision maker. The reading
this supports is the one Section~\ref{sec:geometry} measures: $k$ is the price
the book pays, in log-growth, for exposure to spreads the market has priced as
near-binary.

\section{Price Basis and Execution}
\label{sec:basis}

The price basis is the central empirical question in any spread-selling
backtest. The strategy's edge is a small number obtained by differencing
quoted prices, so a biased basis manufactures or destroys the edge before any
selection logic runs. FCM prices every candidate through three successive
ceilings and four reality gates.

\subsection{Three ceilings on the credit}
\label{sec:ceilings}

\paragraph{1. The smile-fit model credit.} The base credit is
$c_0 = \mathrm{BS}(K_s, \hat\sigma_s) - \mathrm{BS}(K_\ell, \hat\sigma_\ell)$,
Black--Scholes prices at the per-chain fitted volatilities of
Section~\ref{sec:dent}. A fitted surface is used rather than the raw quoted
midpoints of the two legs because a vertical's quoted legs are the two noisiest
points on the chain and their difference is noisier still; the fit pools the
whole chain, both sides, weighted by relative quoted width.

\paragraph{2. The delta ceiling.} A vertical cannot be worth more than the
short leg's own risk-neutral exercise probability times the width, so
\begin{equation}
c_1 \;=\; \min\bigl(c_0,\ |\delta_s|\, w_d\bigr),
\label{eq:deltacap}
\end{equation}
with $\delta_s$ the fitted short-leg delta. This ceiling binds on $5{,}190$ of
$99{,}587$ scored candidates. It exists because the fitted surface can
extrapolate above the arbitrage bound on illiquid strikes: one documented case
booked $0.857$ of a one-dollar width on a book quoted $0.45/2.45$ against
$0.05/0.45$, whose immediately transactable credit was zero.

\paragraph{3. The best-quote ceiling.} The credit cannot exceed the most the
quoted book ever showed:
\begin{equation}
c \;=\; \min\bigl(c_1,\ a_s - \beta_\ell\bigr),
\label{eq:bestcap}
\end{equation}
where $a_s$ is the short leg's ask and $\beta_\ell$ the long leg's bid. This
binds on $2.3\%$ of survivors with a median haircut of $\$0.014$. Both
(\ref{eq:deltacap}) and (\ref{eq:bestcap}) are applied before scoring, so
$\Gamma$ and the daily ordering are computed on the capped credit, not merely
booked on it.

\subsection{Four quote-reality gates}
\label{sec:qualitygates}

Write $\mathrm{touch} = \beta_s - a_\ell$ for the credit obtainable by selling
the short at its bid and buying the long at its ask -- the price available
immediately -- and let $\mu_s, \mu_\ell$ be the two leg midpoints. A candidate
is dropped if any of the following holds.

\begin{enumerate}
\item $\mathrm{touch} \leq 0$: the quoted book cannot pay anything at all.
      $14{,}768$ candidates.
\item $\mu_s < \mu_\ell$: inverted midpoints, the nearer strike quoted below
      the farther one. $625$ candidates.
\item either strike is not a multiple of \$0.50: off-ladder series.
      $499$ candidates.
\item $a_s < \beta_\ell$: quotes crossed across the two strikes.
      $28$ candidates.
\end{enumerate}

Together these remove $15{,}125$ of $114{,}712$ candidates. Gate 3 deserves a
word because its cause is not obvious. A split-adjusted price series keeps
trading on an odd ladder for years after the corporate action -- strikes at
$32.17/32.00$, $196.88/196.25$, $1.67$ increments -- long past the window any
corporate-action gate covers. Those series are thin orphans. Ladder \emph{gaps}
on an ordinary grid, by contrast, are real listings and are kept.

The gates are applied at candidate level, before the daily cap, so the daily
selection backfills from the surviving pool rather than thinning. Candidates
per day fall from $99$ to $84$ and no day falls below the cap of six. Their
cost to the reported book is real and is reported: the in-sample book gives up
$15\%$ of its profit and $0.17$ of its dollar Sharpe ratio. We keep them
because the alternative is a backtest that books spreads no participant could
have transacted.

\subsection{The execution model and the break-even fill}
\label{sec:execution}

Selection and booking are separated. Selection scores the capped model credit
$c$. Booking applies a flat multiple $f$ of that credit, and because $f$ is
flat it changes realized profit without changing any ordering. The reported
case is $f = 1.00$: the trade is booked at exactly the fitted fair value, with
no assumed edge in execution.

Live operation gates on the same number. A candidate qualifies only if its
quoted midpoint is at least $1.00 \times$ the model credit -- below fair value
the spread is being sold for less than the surface says it is worth -- with a
working target of $1.04$ to $1.10\times$.

The sensitivity of the book to $f$ is the single most important fragility in
the system, and Section~\ref{sec:fill} reports the full ladder. By bisection
on realized settlement, total profit crosses zero at
\begin{equation}
f^\ast \;=\; 0.9197 \ \ (\text{in sample}), \qquad
0.8612 \ \ (\text{out of time}), \qquad
0.9123 \ \ (\text{pooled}).
\label{eq:breakeven}
\end{equation}
The strategy requires fills no worse than about $0.92 \times$ the fitted fair
credit. Three independent measurements bear on whether that is met. A sample
of $52$ live trades with a recorded fill price achieved a median of
$1.037 \times$ the model credit, with $47$ of $52$ at or above $1.00\times$.
An earlier sample of $19$ fills averaged $1.08\times$. A sample of $57$ fills
measured against the quoted combination midpoint rather than the model
averaged $0.951\times$ of that midpoint. All three sit above $f^\ast$, and all
three are small samples from a single account; this is the softest quantity in
the paper and we say so again in Section~\ref{sec:discussion}.

\section{Empirical Design}
\label{sec:design}

\subsection{Universe, instrument and candidate construction}

The frame is built from vendor end-of-day option chains on S\&P~100
underlyings plus SPY, QQQ and IWM, January 2020 through September 2026,
$131{,}570$ candidate rows on $93$ names over $1{,}346$ entry days, split
$65{,}255$ bull puts and $66{,}315$ bear calls. On each entry day, Monday
through Thursday, for each name and each side, the system takes the single
vertical expiring the coming Friday -- one to four days to expiry -- whose
short leg has fitted delta nearest $0.55$ within the band $[0.50, 0.60]$, with
the adjacent strike as the long leg, width at most $\$2.50$, and open interest
at least one contract on both legs.

The $0.55$ delta target deserves emphasis, because it determines the character
of the trade and distinguishes it from the out-of-the-money premium sale the
term ``credit spread'' usually denotes. A $0.55$-delta short leg is at or slightly inside the money. The
median selected spread collects $0.497$ of its width, giving payoff odds
$b = 0.989$, and wins outright $47.4\%$ of the time. This is a near-even-money
bet at near-even odds, with the partial zone and the credit cushion supplying
the edge, not a high-probability out-of-the-money sale. The win rate is not a
performance statistic here; it is a design consequence.

\subsection{The gate stack}

Four exclusions apply to both sleeves, each covering entry through expiry plus
one session: scheduled earnings releases, ex-dividend dates, corporate actions
(splits, reverse splits, spinoffs), and -- a separate rule -- any candidate
whose $252$-session $\Preal$ window contains an unadjusted corporate action.
A fifth requirement is historical: a name is not traded until its own close
history holds $252$ completed sessions, which moves the first entry to
2020-09-29.

The earnings calendar is the union of a vendor calendar and dates extracted
from SEC Item 2.02 filings, vetted so that every name-year from 2020 to 2025
carries exactly four reports. The vetting matters: a naive union added $188$
candidate dates of which only $26$ were genuine earnings releases, the rest
being second filings within a quarter -- delivery reports, in-process
research-and-development pre-announcements -- that would have silently deleted
good candidates.

The close store underlying $\Preal$ is split-adjusted before use, with each
corporate action verified against the observed session return: an event is
applied only if the realized move lands within $10\%$ of the implied move and
implies at least a $15\%$ move. Events failing verification are left raw and
handled by the lookback exclusion instead. The asymmetry is deliberate. A
missed adjustment leaves one wrong return in a $252$-point window; a falsely
applied one \emph{inserts} a spurious jump that was not there, which is worse.
A loose tolerance in an earlier revision falsely verified a $4.9\%$ spinoff
factor against an ordinary $3.7\%$ session and rescaled an entire price
prefix.

\subsection{Direction}

Bull puts fire every entry day. Bear calls fire only when the prior completed
SPY close is below its $100$-session simple moving average, and additionally
require fitted short-leg implied volatility below $0.35$. The asymmetry is
empirical and Section~\ref{sec:gates} measures it: bear calls taken every day
are strongly negative, bear calls restricted to the regime are the book's
drawdown insurance.

\subsection{Selection}
\label{sec:selection}

Both sleeves pass a credit floor,
\begin{equation}
c / w_d \;\geq\; 0.45 ,
\label{eq:floor}
\end{equation}
and then compete in a single pooled ranking: the six highest-$\Gamma$
candidates per entry day are taken, both sides together. There is no
$\Gamma$ threshold.

The funnel is as follows. Of $131{,}570$ frame rows, $93{,}450$ survive the
regime assignment, the four gates, the four quote-reality gates and the
$252$-session requirement. Of those, $59{,}819$ clear the credit floor, and
$25{,}228$ are eligible once the bear sleeve's regime and volatility
conditions are imposed -- $22{,}857$ bull puts and $2{,}371$ bear calls, on
$92$ names across $1{,}120$ entry days, a mean of $22.5$ eligible candidates
per day. The daily cap of six selects $6{,}593$ of them.

\subsection{Why there is no threshold}
\label{sec:nothreshold}

Proposition~\ref{prop:halfplane} invites an acceptance threshold
$\Gamma \geq \theta$, and earlier configurations of this system carried one.
FCM has none, and the reason is the drift-free $\Preal$ of
Section~\ref{sec:preal}. Removing the window drift shrank per-candidate edges by
roughly an order of magnitude -- the median selected $\Gamma$ is
$1.21\times10^{-4}$ in sample -- so a threshold calibrated on raw-drift edges
admits almost nothing, and one recalibrated downward admits almost everything.
With a cap of six against a pool of $22.5$, the cap is the binding constraint
and the threshold is redundant.

The consequence is stated plainly: negative-$\Gamma$ candidates are selected
when fewer than six positive ones clear the floor. In sample this happens on
$992$ of $5{,}779$ picks ($17.2\%$), which together earned $+\$5{,}502$ per
contract; out of time it happens on $54$ of $814$, earning $+\$1{,}262$.
A threshold would have excluded profitable trades.

\subsection{Settlement and sizing}
\label{sec:sizing}

Settlement uses the underlying's expiry-day close. A full win pays the booked
credit; a full loss pays the booked maximum loss; a partial outcome pays
$c - |S_T - K_s|$, and a \emph{profitable} partial is haircut to $50\%$ of its
value to reflect pin and assignment risk, while partial losses are taken in
full. Commissions are set to zero throughout.

The reporting arm is a flat two contracts per pick on a \$10{,}000 starting
bankroll. Three alternative arms are reported: flat one contract, and
fractional-Kelly sizing at one quarter and one half of
$w^\ast_\kappa$ from (\ref{eq:kellycarry}), each capped at five contracts, in
the tradition of \cite{maclean2011}. Candidates for which carry-adjusted Kelly
declines to size -- exactly the growth-negative rows -- are sized at one
contract, not at the floor stake; an earlier version fell back to the
uncarried $w^\ast$ for those rows, which clipped to $0.01$ and, on a small
maximum loss, sized \emph{above} one contract, resurrecting trades Kelly had
rejected.

\section{Results}
\label{sec:results}

Unless stated otherwise every figure is at $f = 1.00 \times$ model credit,
two contracts per pick, \$10{,}000 starting bankroll, and the Sharpe ratio
quoted is the annualized dollar-P\&L Sharpe ratio, which we prefer to the
return-based ratio because a flat-contract book has no well-defined invested
base.

\subsection{Headline}
\label{sec:headline}

\begin{table}[h]
\centering
\caption{FCM configuration, $k = 24$, pooled top six per day, $f = 1.00\times$
model credit, \$10{,}000 starting bankroll.}
\label{tab:headline}
\begin{tabular}{lrr}
\toprule
 & In sample & Out of time \\
 & 2020-09-29 to 2026-01-02 & 2026-01-05 to 2026-09-25 \\
\midrule
Trades & 5{,}779 & 814 \\
Trading days & 1{,}322 & 183 \\
Entry days used & 984 & 136 \\
Names traded & 90 & 73 \\
Final equity, 2 contracts & \$83{,}391 & \$28{,}198 \\
Total return & $+733.9\%$ & $+182.0\%$ \\
CAGR & $49.7\%$ & --- \\
Sharpe (daily \$, annualized) & $+1.26$ & $+2.86$ \\
Sharpe (weekly, annualized) & $+1.29$ & $+3.19$ \\
Maximum drawdown & $-30.3\%$ & $-15.2\%$ \\
Yield on dollars risked & $7.8\%$ & $13.6\%$ \\
Final equity, 1 contract & \$46{,}695 & \$19{,}099 \\
Final equity, $\tfrac14$ Kelly (cap 5) & \$91{,}828 & \$31{,}502 \\
Final equity, $\tfrac12$ Kelly (cap 5) & \$128{,}198 & \$42{,}180 \\
SPY over the same window & \$22{,}113 & \$11{,}276 \\
SPY Sharpe / drawdown & $+0.97$ / $-24.5\%$ & $+1.39$ / $-8.9\%$ \\
\midrule
Win / partial / loss & 2{,}741 / 851 / 2{,}187 & 391 / 143 / 280 \\
Win rate & $47.4\%$ & $48.0\%$ \\
Mean profit per contract & \$6.35 & \$11.18 \\
Median payoff odds $b$ & 0.989 & 1.017 \\
Median credit / width & 0.497 & 0.504 \\
Median $\dent$ & 0.1001 & 0.0962 \\
Median $E$ & 0.00151 & 0.00332 \\
Median $\Gamma$ & $1.21\times10^{-4}$ & $2.75\times10^{-4}$ \\
Median $|\delta_s|$ & 0.558 & 0.553 \\
\bottomrule
\end{tabular}
\end{table}

The gross flows are worth recording because they set the scale of everything
that follows. In sample, at one contract, full wins paid $+\$223{,}949$, full
losses cost $-\$178{,}288$, and partial outcomes cost $-\$8{,}965$, netting
$+\$36{,}695$ on $\$411{,}202$ of gross settlement. The book is a $9\%$
residual of two large opposing flows. Nothing in this system is a large per
trade effect, and any statement about it that cannot survive a few percent of
error in the credit basis is not a statement about the strategy.

\subsection{Year by year}

\begin{table}[h]
\centering
\caption{By entry year, one contract per pick, $f = 1.00\times$ model.
``Yield'' is profit over dollars at risk.}
\label{tab:byyear}
\begin{tabular}{lrrrrrrr}
\toprule
Year & Trades & Profit & Per trade & Win & Loss & Yield & Bull / bear \\
\midrule
2020 (part) & 16    & \$275      & \$17.18  & $50.0\%$ & $37.5\%$ & $13.6\%$ & 16 / 0 \\
2021        & 1{,}142 & \$11{,}359 & \$9.95   & $51.1\%$ & $35.6\%$ & $12.6\%$ & 1{,}133 / 9 \\
2022        & 1{,}172 & $-\$749$   & $-\$0.64$& $41.2\%$ & $44.8\%$ & $-0.8\%$ & 775 / 397 \\
2023        & 1{,}151 & \$6{,}361  & \$5.53   & $48.1\%$ & $38.8\%$ & $6.8\%$  & 980 / 171 \\
2024        & 1{,}151 & \$10{,}901 & \$9.47   & $49.0\%$ & $34.8\%$ & $11.4\%$ & 1{,}143 / 8 \\
2025        & 1{,}147 & \$8{,}549  & \$7.45   & $48.0\%$ & $35.1\%$ & $8.8\%$  & 1{,}057 / 90 \\
2026 (OOT)  & 814     & \$9{,}099  & \$11.18  & $48.0\%$ & $34.4\%$ & $13.6\%$ & 770 / 44 \\
\bottomrule
\end{tabular}
\end{table}

One year is negative: 2022, at $-\$749$ on $1{,}172$ trades, essentially flat.
It is also the only year in which the bear sleeve was large, $397$ of
$1{,}172$ picks, because it is the only year SPY spent substantial time below
its $100$-session average. The sleeve did not save 2022, but
Section~\ref{sec:gates} shows what the book looks like without it.

\subsection{What the ranking contributes}
\label{sec:ranker}

As in the first three papers of this series, the benchmark is the Kelly
criterion. Because FCM applies no threshold, the question is not whether a gate
earns its place but whether the entropy-discounted ordering beats the pure
growth ordering. The comparison that isolates it holds the eligible pool, every
gate, the cap and the fill identical and varies only the sort key; ranking on
$E$ alone is precisely Kelly criterion selection under the trader's measure.

\begin{table}[h]
\centering
\caption{Ranker ablation. Identical eligible pool ($25{,}228$ candidates),
identical gates, pooled top six per day, $f = 1.00\times$ model.}
\label{tab:ranker}
\begin{tabular}{lrrrrrr}
\toprule
 & \multicolumn{3}{c}{In sample} & \multicolumn{3}{c}{Out of time} \\
\cmidrule(lr){2-4}\cmidrule(lr){5-7}
Sort key & Profit & Sharpe & Max DD & Profit & Sharpe & Max DD \\
\midrule
$\Gamma = E\,e^{-24\dent}$ (FCM) & \$73{,}391 & $+1.26$ & $-30.3\%$ & \$18{,}198 & $+2.86$ & $-15.2\%$ \\
$E$ alone (Kelly growth, $k=0$)  & \$69{,}401 & $+1.13$ & $-50.1\%$ & \$16{,}464 & $+2.73$ & $-10.2\%$ \\
$-\dent$ alone (max market entropy) & \$51{,}569 & $+0.90$ & $-72.4\%$ & \$9{,}669 & $+1.44$ & $-17.2\%$ \\
Credit / width                    & \$67{,}035 & $+1.06$ & $-54.9\%$ & \$8{,}668 & $+1.31$ & $-35.4\%$ \\
Random                            & \$48{,}973 & $+0.77$ & $-122.9\%$ & \$917 & $+0.14$ & $-64.0\%$ \\
$\Gamma$ ascending (worst six)    & \$8{,}105 & $+0.13$ & $-149.5\%$ & $-\$8{,}362$ & $-1.05$ & $-128.9\%$ \\
\bottomrule
\end{tabular}
\end{table}

Three readings. First, the ordering is correct in both windows and in the
right direction at both ends: $\Gamma$ beats each of its own factors taken
alone, each factor beats random, and reversing $\Gamma$ nearly destroys the
book and turns the out-of-time window negative. The score carries information.

Second, $\Gamma$'s advantage over the growth factor alone is modest in dollars
-- $+6\%$ in sample, $+11\%$ out of time -- and large in path: maximum
drawdown falls from $-50.1\%$ to $-30.3\%$ in sample. The entropy factor's
contribution is risk control, not return. This is a more modest claim than a
reader of the first three papers might expect, and it is the claim the data
support.

Third, the eligible pool is profitable before any ranking: taking all
$25{,}228$ eligible candidates at one contract would have earned
$+\$80{,}396$, or \$3.19 per trade, against the selected book's \$6.94 per
trade on $6{,}593$ trades. The ranker roughly doubles the per-trade rate. The
gates and the credit floor, not the score, are what make the pool positive in
the first place, and Section~\ref{sec:gates} measures them.

\subsection{The entropy strength \texorpdfstring{$k$}{k}}
\label{sec:ksweep}

\begin{table}[h]
\centering
\caption{$k$ sweep on $\Gamma$-rank. Same pool and cap; only the exponent
changes. $f = 1.00\times$ model.}
\label{tab:ksweep}
\begin{tabular}{rrrrrrr}
\toprule
 & \multicolumn{3}{c}{In sample} & \multicolumn{3}{c}{Out of time} \\
\cmidrule(lr){2-4}\cmidrule(lr){5-7}
$k$ & Profit & Sharpe & Max DD & Profit & Sharpe & Max DD \\
\midrule
0  & \$69{,}401 & $+1.13$ & $-50.1\%$ & \$16{,}464 & $+2.73$ & $-10.2\%$ \\
2  & \$65{,}398 & $+1.06$ & $-52.7\%$ & \$16{,}698 & $+2.81$ & $-10.1\%$ \\
4  & \$65{,}443 & $+1.06$ & $-49.5\%$ & \$17{,}396 & $+2.98$ & $-9.4\%$ \\
8  & \$69{,}484 & $+1.14$ & $-39.3\%$ & \$16{,}472 & $+2.68$ & $-13.4\%$ \\
12 & \$70{,}021 & $+1.17$ & $-39.4\%$ & \$16{,}422 & $+2.65$ & $-14.5\%$ \\
16 & \$70{,}867 & $+1.20$ & $-37.7\%$ & \$17{,}317 & $+2.67$ & $-15.3\%$ \\
20 & \$71{,}815 & $+1.23$ & $-32.7\%$ & \$18{,}118 & $+2.88$ & $-14.9\%$ \\
\textbf{24} & \textbf{\$73{,}391} & $\mathbf{+1.26}$ & $\mathbf{-30.3\%}$ & \textbf{\$18{,}198} & $\mathbf{+2.86}$ & $\mathbf{-15.2\%}$ \\
28 & \$75{,}570 & $+1.30$ & $-29.8\%$ & \$18{,}548 & $+2.96$ & $-14.8\%$ \\
32 & \$71{,}819 & $+1.24$ & $-31.7\%$ & \$18{,}588 & $+3.02$ & $-14.6\%$ \\
48 & \$65{,}001 & $+1.12$ & $-32.4\%$ & \$16{,}844 & $+2.62$ & $-15.7\%$ \\
64 & \$67{,}546 & $+1.16$ & $-33.2\%$ & \$17{,}600 & $+2.79$ & $-14.1\%$ \\
\bottomrule
\end{tabular}
\end{table}

The surface has a clear shape and one honest wrinkle. In sample the Sharpe
ratio is $1.13$ at $k = 0$, dips slightly to $1.06$ near $k = 2$--$4$, then
rises monotonically to a peak of $1.30$ at $k = 28$ before falling away by
$k = 48$; maximum drawdown improves almost monotonically from $-50\%$ to
$-30\%$ over the same range. Out of time the surface is nearly flat, $2.6$ to
$3.0$ throughout, with no interior peak worth naming. The plateau
$k \in [20, 32]$ is within $3\%$ of the in-sample peak.

The wrinkle: the canonical $k = 24$ is \emph{not} the in-sample maximum.
$k = 28$ is better on both windows on this pool. The value $24$ was fixed on
an earlier sweep run before the quote-reality gates of
Section~\ref{sec:qualitygates} and the earnings-source correction were
applied, and we have not moved it, because moving a parameter onto the peak of
whichever pool was most recently rebuilt is how a configuration becomes
fitted. What the sweep supports is the plateau, not the point: the claim this
paper makes is that $k$ in the low twenties to low thirties is right, and that
the quantity it buys is drawdown.

\subsection{What each gate is worth}
\label{sec:gates}

\begin{table}[h]
\centering
\caption{Gate ablation. $\Gamma$-rank, pooled top six per day,
$f = 1.00\times$ model. One condition changed per row.}
\label{tab:gates}
\begin{tabular}{lrrrrrr}
\toprule
 & \multicolumn{3}{c}{In sample} & \multicolumn{3}{c}{Out of time} \\
\cmidrule(lr){2-4}\cmidrule(lr){5-7}
Configuration & Profit & Sharpe & Max DD & Profit & Sharpe & Max DD \\
\midrule
FCM as specified              & \$73{,}391 & $+1.26$ & $-30.3\%$ & \$18{,}198 & $+2.86$ & $-15.2\%$ \\
\midrule
No credit floor               & \$68{,}890 & $+1.20$ & $-32.5\%$ & \$17{,}112 & $+2.62$ & $-17.2\%$ \\
Floor $0.40$ not $0.45$       & \$65{,}519 & $+1.10$ & $-39.8\%$ & \$17{,}848 & $+2.70$ & $-16.9\%$ \\
Floor $0.50$ not $0.45$       & \$48{,}151 & $+0.86$ & $-59.2\%$ & \$15{,}657 & $+2.52$ & $-12.9\%$ \\
\midrule
No ex-dividend gate           & \$42{,}185 & $+0.76$ & $-42.9\%$ & \$14{,}256 & $+2.23$ & $-18.5\%$ \\
No earnings gate              & \$61{,}568 & $+1.09$ & $-31.8\%$ & \$17{,}741 & $+2.95$ & $-11.4\%$ \\
No corporate-action gates     & \$75{,}004 & $+1.30$ & $-33.7\%$ & \$19{,}936 & $+3.13$ & $-14.8\%$ \\
\midrule
Bears every day (no regime)   & \$22{,}170 & $+0.55$ & $-86.0\%$ & \$9{,}475  & $+2.16$ & $-28.0\%$ \\
No bear sleeve at all         & \$57{,}654 & $+0.86$ & $-58.9\%$ & \$9{,}255  & $+1.37$ & $-45.4\%$ \\
Bears without the IV $<0.35$ cut & \$65{,}315 & $+1.08$ & $-40.1\%$ & \$22{,}712 & $+3.42$ & $-12.5\%$ \\
\bottomrule
\end{tabular}
\end{table}

The credit floor is a genuine optimum rather than a convenience: $0.45$ beats
no floor, beats $0.40$, and beats $0.50$ decisively. The economic content is
that credit over width is the market's own full-loss odds, so a floor on it is
a floor on the compensation per unit of tail exposure, and raising it too far
selects only the spreads the market has already marked as likely to lose.

The ex-dividend gate is the single most valuable exclusion in the system,
worth $\$31{,}206$ and $0.50$ of Sharpe ratio in sample. The earnings gate is
worth $\$11{,}823$ and $0.17$. The corporate-action gates \emph{cost}
$\$1{,}613$ and $0.04$ in sample and more out of time, and we keep them
anyway, which is the clearest instance in this paper of a correctness choice
taken against performance: a book that holds a position across a spinoff is
booking a price that did not exist.

The bear sleeve is the result that most changes how the system should be
understood. It is small -- $2{,}371$ of $25{,}228$ eligible candidates, $719$
of $6{,}593$ picks, $10.9\%$ -- and removing it costs far more than its size:
in sample $\$73{,}391 \to \$57{,}654$ with drawdown deepening from $-30.3\%$
to $-58.9\%$, and out of time $\$18{,}198 \to \$9{,}255$ with drawdown from
$-15.2\%$ to $-45.4\%$. The sleeve is the book's hedge, concentrated in the
periods when the bull sleeve is being punished. Removing its regime
restriction, on the other hand, is catastrophic in sample
($\$22{,}170$, Sharpe $0.55$, drawdown $-86\%$): bear calls sold
indiscriminately in a rising market are the wrong side of the same variance
premium. The volatility cut on the sleeve is the one component with a genuine
in-sample / out-of-time disagreement -- removing it costs $\$8{,}076$ in
sample but gains $\$4{,}514$ out of time -- and is flagged as unresolved.

\subsection{The daily cap}
\label{sec:cap}

\begin{table}[h]
\centering
\caption{Daily cap on $\Gamma$-rank, $f = 1.00\times$ model.}
\label{tab:cap}
\begin{tabular}{rrrrrrrr}
\toprule
 & \multicolumn{3}{c}{In sample} & \multicolumn{4}{c}{Out of time} \\
\cmidrule(lr){2-4}\cmidrule(lr){5-8}
Cap & Trades & Profit & Sharpe & Trades & Profit & Sharpe & Max DD \\
\midrule
1  & 984     & \$13{,}439  & $+0.97$ & 136     & \$7{,}533  & $+5.29$ & $-2.3\%$ \\
2  & 1{,}958 & \$14{,}997  & $+0.65$ & 272     & \$8{,}430  & $+3.14$ & $-8.8\%$ \\
3  & 2{,}927 & \$34{,}259  & $+1.09$ & 408     & \$10{,}723 & $+2.85$ & $-12.2\%$ \\
4  & 3{,}892 & \$50{,}273  & $+1.22$ & 544     & \$13{,}490 & $+2.70$ & $-12.7\%$ \\
\textbf{6} & \textbf{5{,}779} & \textbf{\$73{,}391} & $\mathbf{+1.26}$ & \textbf{814} & \textbf{\$18{,}198} & $\mathbf{+2.86}$ & $\mathbf{-15.2\%}$ \\
8  & 7{,}562 & \$82{,}444  & $+1.06$ & 1{,}084 & \$19{,}828 & $+2.74$ & $-10.4\%$ \\
10 & 9{,}251 & \$102{,}410 & $+1.07$ & 1{,}346 & \$26{,}995 & $+3.14$ & $-8.5\%$ \\
15 & 12{,}970& \$127{,}076 & $+0.95$ & 1{,}948 & \$24{,}811 & $+1.97$ & $-33.3\%$ \\
\bottomrule
\end{tabular}
\end{table}

Six is the in-sample Sharpe maximum and the cap in production. Dollars keep
rising to a cap of fifteen, so the choice of six is a risk-adjusted choice
and nothing else. The small-cap rows carry a warning about reading
out-of-time Sharpe ratios on this instrument: a cap of one produces an
out-of-time Sharpe ratio of $5.29$ on $136$ trades and $7{,}533$ dollars, which
is a statement about $136$ Friday closes and not about a strategy.

\subsection{The fill ladder}
\label{sec:fill}

\begin{table}[h]
\centering
\caption{Fill sensitivity. Selection is identical in every row; only the
booked fraction of the model credit changes.}
\label{tab:fill}
\begin{tabular}{lrrrrr}
\toprule
 & \multicolumn{2}{c}{In sample} & \multicolumn{3}{c}{Out of time} \\
\cmidrule(lr){2-3}\cmidrule(lr){4-6}
Fill $f$ & Profit & Sharpe & Profit & Sharpe & Max DD \\
\midrule
$0.80\times$ model & $-\$109{,}731$ & $-1.88$ & $-\$8{,}050$ & $-1.26$ & $-82.1\%$ \\
$0.85\times$       & $-\$63{,}823$  & $-1.10$ & $-\$1{,}471$ & $-0.23$ & $-41.0\%$ \\
$0.90\times$       & $-\$17{,}998$  & $-0.31$ & \$5{,}094    & $+0.80$ & $-23.3\%$ \\
$0.95\times$       & \$27{,}724     & $+0.48$ & \$11{,}651   & $+1.84$ & $-18.6\%$ \\
$\mathbf{1.00\times}$ & \textbf{\$73{,}391} & $\mathbf{+1.26}$ & \textbf{\$18{,}198} & $\mathbf{+2.86}$ & $\mathbf{-15.2\%}$ \\
$1.04\times$       & \$109{,}849    & $+1.89$ & \$23{,}424   & $+3.68$ & $-13.1\%$ \\
$1.10\times$       & \$164{,}430    & $+2.81$ & \$31{,}244   & $+4.87$ & $-10.5\%$ \\
\bottomrule
\end{tabular}
\end{table}

Each five percentage points of fill quality is worth roughly \$45{,}000 of
in-sample profit at two contracts, and the sign of the whole enterprise turns
between $0.90$ and $0.95$. The exact crossings are at (\ref{eq:breakeven}):
$0.9197\times$ in sample, $0.8612\times$ out of time. The reported case sits
$8$ points above the in-sample crossing and the measured fills sit $12$ points
above it, which is a margin but not a comfortable one. Readers should treat
every other number in this paper as conditional on execution at or above the
fitted fair value, and Section~\ref{sec:live} reports what happened when the
criterion was run against live quotes.

\subsection{A worked example}
\label{sec:example}

Consider the bull put credit spread on USB selected on Wednesday 2026-01-28
for the Friday 2026-01-30 expiry, two days to expiry, with the underlying at
\$55.74.

The short leg is the \$56 put, quoted $0.44/0.64$; the long leg is the \$55
put, quoted $0.13/0.21$. Note the short strike is \emph{above} spot: at a
$0.55$ delta target the short leg is slightly in the money, with fitted delta
$|\delta_s| = 0.5815$. The chain smile fits implied volatilities of $0.2906$
at the short strike and $0.2601$ at the long, giving a Black--Scholes credit of
$c_0 = 0.4667$ on a \$1.00 width. The delta ceiling (\ref{eq:deltacap}) is
$0.5815 \times 1.00 = 0.5815$ and the best-quote ceiling (\ref{eq:bestcap}) is
$0.64 - 0.13 = 0.51$; neither binds, so $c = 0.4667$. The immediately
transactable credit, $\mathrm{touch} = 0.44 - 0.21 = 0.23$, is positive, so
gate 1 passes. Maximum loss is $m = 0.5333$ and the payoff odds are
$b = 0.8751$, so $\alpha = (b-1)/(2b) = -0.0714$ and the partial state pays
$\alpha b = -0.0625$ per unit risked.

The market measure from (\ref{eq:nd2})--(\ref{eq:qbs}) is
$\Qbs = (0.4102,\ 0.2468,\ 0.3430)$. Its Shannon entropy is $H = 1.0779$, so
\begin{equation*}
\dent \;=\; \ln 3 - H \;=\; 1.0986 - 1.0779 \;=\; 0.0207 ,
\end{equation*}
a low value: the market assigns $34.3\%$ to the partial zone, a near-uniform
three-state view, exactly the high-entropy situation the criterion is built to
prefer. Counting the $252$ trailing two-session moves of USB against the two
strikes, demeaned, gives
$\Preal = (p, q, r_0) = (0.3965,\ 0.2229,\ 0.3807)$: the realized history
assigns \emph{less} probability to the full loss than the market does
($0.223$ against $0.247$) and more to the partial zone.

The first-order condition of (\ref{eq:kelly3}) gives $w^\ast = 0.1869$ and
\begin{align*}
\ell(w^\ast) &= 0.3965\ln(1.1636) + 0.3807\ln(0.9883) + 0.2229\ln(0.8131)\\
&= 0.0600 - 0.0045 - 0.0461 \;=\; 0.00947 ,
\end{align*}
so $E = e^{0.00947} - 1 = +0.952\%$ and
\begin{equation*}
\Gamma \;=\; 0.009518 \times e^{-24 \times 0.020728}
\;=\; 0.009518 \times 0.6081 \;=\; 0.005788 ,
\end{equation*}
which is in the top decile of its day and the spread was selected. USB closed
on the Friday at \$56.11, above the short strike, and the spread settled a
full win paying the credit, $+\$46.67$ per contract.

The example also illustrates the scale of the edge. The Kelly expected value
is under one percent, the entropy discount multiplies it by $0.61$, and the
resulting score is $5.8 \times 10^{-3}$ — a strong candidate in a pool whose
median is $1.2\times10^{-4}$. These are small numbers by construction.

\subsection{What \texorpdfstring{$\Gamma$}{Gamma} ranks, and what it does not}
\label{sec:concordance}

This is the finding that most constrains how the criterion should be used, and
it is negative.

\begin{table}[h]
\centering
\caption{Realized outcomes by $\Gamma$ decile within the selected book,
$f = 1.00\times$ model, one contract. Decile 1 is lowest $\Gamma$.}
\label{tab:decile}
\begin{tabular}{rrrrrrr}
\toprule
 & \multicolumn{3}{c}{In sample ($n=5{,}779$)} & \multicolumn{3}{c}{Out of time ($n=814$)} \\
\cmidrule(lr){2-4}\cmidrule(lr){5-7}
Decile & Full-loss rate & \$ / trade & median $\dent$ & Full-loss rate & \$ / trade & median $\dent$ \\
\midrule
1  & $39.6\%$ & \$6.14  & 0.134 & $31.7\%$ & \$22.90  & 0.211 \\
2  & $42.6\%$ & \$5.22  & 0.218 & $45.7\%$ & \$1.08   & 0.146 \\
3  & $43.8\%$ & $-\$2.21$ & 0.144 & $38.3\%$ & \$2.85  & 0.119 \\
4  & $42.0\%$ & \$8.30  & 0.134 & $28.0\%$ & \$21.40  & 0.118 \\
5  & $39.8\%$ & \$3.11  & 0.109 & $33.3\%$ & \$11.22  & 0.110 \\
6  & $35.5\%$ & \$9.00  & 0.100 & $45.7\%$ & $-\$0.99$ & 0.097 \\
7  & $37.4\%$ & \$6.49  & 0.091 & $39.0\%$ & \$11.25  & 0.098 \\
8  & $35.5\%$ & \$11.39 & 0.076 & $30.9\%$ & \$11.49  & 0.069 \\
9  & $32.4\%$ & \$8.18  & 0.061 & $28.4\%$ & \$13.11  & 0.052 \\
10 & $29.9\%$ & \$7.89  & 0.038 & $23.2\%$ & \$17.13  & 0.036 \\
\bottomrule
\end{tabular}
\end{table}

The full-loss rate falls across deciles, from $39.6\%$ to $29.9\%$ in sample
and from the low forties to $23.2\%$ out of time, and the Spearman correlation
between $\Gamma$ and a full-loss indicator is $-0.079$ in sample and $-0.070$
out of time. Dollars per trade show no such pattern: the Spearman correlation
between $\Gamma$ and realized profit is $+0.013$ in sample and $-0.008$ out of
time, statistically indistinguishable from zero in both.

$\Gamma$ orders the probability of losing. It does not order dollars.

The mechanism is visible in Table~\ref{tab:decile}'s third column and in the
geometry of Section~\ref{sec:geometry}. High $\Gamma$ means low $\dent$ means
a fat partial zone, and a fat partial zone converts would-be full losses into
partial outcomes. But a fat partial zone also means the market's expected move
is small relative to the width, which means a smaller credit for the same
width. The criterion buys a lower loss rate and pays for it in credit size,
and on this instrument family the two effects very nearly cancel in dollars.

Table~\ref{tab:conc} measures the same thing on the whole eligible pool with a
cleaner statistic: within-day pairwise concordance, the probability that of two
candidates seen on the same day the higher-scored one realizes the better
outcome. Chance is $0.500$.

\begin{table}[h]
\centering
\caption{Within-day pairwise concordance on the eligible pool.
``\$'' target is realized profit; ``not a loss'' target is avoidance of the
full-loss state. $307{,}389$ in-sample and $49{,}331$ out-of-time pairs.}
\label{tab:conc}
\begin{tabular}{lrrrr}
\toprule
 & \multicolumn{2}{c}{In sample} & \multicolumn{2}{c}{Out of time} \\
\cmidrule(lr){2-3}\cmidrule(lr){4-5}
Score & \$ target & not a loss & \$ target & not a loss \\
\midrule
$\Gamma = E\,e^{-24\dent}$ & $0.5225$ & $0.5163$ & $0.5348$ & $0.5497$ \\
$E$ alone                  & $0.5123$ & $0.5255$ & $0.5248$ & $0.5541$ \\
$-\dent$ alone             & $\mathit{0.4801}$ & $\mathbf{0.5596}$ & $\mathit{0.4921}$ & $0.5560$ \\
Credit / width             & $\mathbf{0.5821}$ & $\mathit{0.4509}$ & $\mathbf{0.5947}$ & $\mathit{0.4872}$ \\
\bottomrule
\end{tabular}
\end{table}

The pattern is the paper's cleanest theoretical result, measured. The entropy
factor alone is the best of the four at predicting which spreads avoid a full
loss ($0.560$ in sample) and is \emph{below chance} at predicting dollars
($0.480$). Credit over width is the best at predicting dollars ($0.582$) --
mechanically so, since for a given settlement a richer credit books more -- and
is \emph{below chance} at predicting loss avoidance ($0.451$), again
mechanically, since a richer credit on a fixed width is the market's own
statement that a loss is more likely. The two targets are in direct tension,
and the two factors sit on opposite sides of it. Their product, $\Gamma$, is
the only one of the four scores tested that is above chance on both targets in
both windows. That is what the two-factor form is for, and it is also precisely
why $\Gamma$ cannot be expected to maximize dollars: it is buying a position on
both axes rather than a corner on one.

The practical consequence, which Section~\ref{sec:discussion} returns to, is
that per-trade discrimination is thin everywhere. A concordance of $0.52$ is a
genuine but small effect. The book's advantage over random selection
(Table~\ref{tab:ranker}: Sharpe $1.26$ against $0.77$) is produced at the
extremes of the distribution and by the credit floor, not by a strong
per-trade classifier.

\subsection{Diversification within the day}
\label{sec:diversification}

The six daily positions are not independent. Pooling within-day demeaned
cross-products over $1{,}110$ entry days gives an average pairwise correlation
of realized profit of $0.161$, so at the realized mean cap of $5.93$
positions the book behaves like
\begin{equation}
n_{\mathrm{eff}} \;=\; \frac{n}{1 + (n-1)\bar\rho} \;=\;
\frac{5.93}{1 + 4.93 \times 0.161} \;=\; 3.30
\label{eq:neff}
\end{equation}
independent bets per day rather than six. The cause is concentration on one
side: $88.6\%$ of picks are bull puts, so the daily book is a partly
undiversified long-equity position.

That exposure is measurable and asymmetric in time. Correlation between the
book's entry-day profit and the same day's SPY return is $-0.011$ on
$1{,}120$ days -- the entry decision carries no same-day market bet.
Correlation between settlement-day profit and the same day's SPY return is
$+0.444$ on $299$ settlement days. The strategy is market-neutral at entry and
long beta at expiry, which is the correct description of a book of mostly
short puts, and it is the reason the regime-gated bear sleeve does as much work
in Table~\ref{tab:gates} as its small size would not suggest.

\subsection{Reproducibility across data sources}
\label{sec:reproducibility}

The claim in Section~\ref{sec:dent} that $\dent$ is the robust factor and the
growth factor the fragile one rests on a reconciliation between the vendor
end-of-day frame used above and independent brokerage snapshots taken at
15:30 Eastern, over twenty shared entry days in August and September 2026.
On $662$ shared name-days, of which $463$ carry identical strikes on both
sides, the components reproduce as follows: $\dent$ at rank correlation
$0.98$, the smile-fit model credit at $0.974$, the Kelly expected value at
$0.41$, and the estimated full-win probability at $0.54$. The composite score
reproduces at a within-day Spearman median of $0.49$--$0.52$ and a pairwise
concordance of $0.692$ against $0.500$ for unrelated.

The cause of the win-probability gap was isolated and is instructive:
recomputing $\Preal$ on the brokerage snapshot's 15:30 spot rather than the
vendor's closing spot raises the correlation from $0.53$ to $0.90$ and drops
the median absolute difference from $0.024$ to $0.000$. The two spots differ
by only $0.14\%$ at the median, but $\Preal$ at a near-the-money strike is a
knife edge: a $0.3\%$ spot move shifts $p$ by $0.043$ at the median. The
growth factor is not merely noisy in estimation; it is structurally sensitive
to the entry spot in a way the entropy factor is not. A $0.55$-delta design
pays for its favourable payoff odds with exactly this sensitivity.

The honest reading of these numbers is that the two sources agree far better
on what to \emph{avoid} than on what to take: $83\%$ of the brokerage data's
worst-quintile candidates fall in the vendor data's bottom half, against
$64\%$ agreement at the top. A top-six overlap statistic between the two
sources is $18$--$38\%$ depending on which rules are aligned, and is the wrong
yardstick; quintile agreement and the bottom-half exclusion are the
informative ones.

\subsection{Live quotes}
\label{sec:live}

The criterion has been run against archived live brokerage snapshots, which is
the only evaluation in this paper that uses neither vendor quotes nor
end-of-day timing. Every 15:30 and 15:45 snapshot from 2026-08-20 to
2026-10-02 -- $29$ scan days, $84$ names, $1{,}893$ candidate rows -- was
rescored under the configuration of Section~\ref{sec:design}: credit floor
$0.45$, quoted midpoint at least $1.00\times$ model credit, pooled top six by
$\Gamma$ at $k = 24$, with the 15:45 scan topping the day up to six. SPY sat
above its $100$-session average throughout the window, so the book is bull
puts only.

The result is $146$ trades: $66$ full wins, $22$ partials, $58$ full losses, a
$45.2\%$ win rate, $+\$2{,}515$ at two contracts and $+\$1{,}258$ at one
(\$8.61 per contract), an annualized dollar Sharpe ratio of $+3.54$, a maximum
drawdown of $-6.0\%$ and a yield on dollars risked of $9.5\%$. Booking is at
$1.00\times$ model credit, the same assumption as the vendor book.

Two caveats are necessary and neither is small. Six weeks and $146$ trades
cannot distinguish a Sharpe ratio of $3.5$ from one of $1.0$; the window's
value is that it shows the criterion is computable and positive on quotes a
participant could actually see, not that it confirms the level. And the vendor
book and the live replay disagree on \emph{which} names to take over these
same weeks, sharing roughly $30\%$ of picks, for the reasons decomposed in
Section~\ref{sec:reproducibility}: a quote gate on a different quote source,
and a candidate universe built from a 15:30 chain rather than a closing one.
Both books are positive over the window; they are not the same book.

\section{Discussion}
\label{sec:discussion}

\subsection{What the entropy factor contributes}

Taking the penalty as $\dkl(\Qbs \| U_3)$ gives the construction three
properties. The entropy factor is a measurable property of the instrument
rather than of a forecast (Section~\ref{sec:geometry}), with an interpretation
-- is this a genuine three-state wager or a disguised binary one -- that can be
checked on any chain without estimating anything. It reproduces across
independent data sources at a rank correlation of $0.98$, so the score's only
fragile input is the growth factor and the fragility is confined there
(Section~\ref{sec:reproducibility}). And it is the published GROUND form of
\cite{mercurio2020binary, mercurio2020gepo}, so the convex-risk-measure result
of that paper applies to the penalty as stated rather than by analogy. What it
does not give is a robust-control reading: the penalty is a state of the market,
not a charge on the agent's own beliefs (Section~\ref{sec:hs}).

Against that, the entropy factor's contribution to the book is smaller than the
emphasis this paper places on it might suggest. Growth selection alone earns
$\$69{,}401$ against $\Gamma$'s $\$73{,}391$, and the entropy factor's real
contribution is to cut maximum drawdown from $-50\%$ to $-30\%$. That is a
worthwhile contribution and a modest claim, and the modest claim is the one the
data support.

\subsection{The criterion ranks loss rate, not dollars}

Section~\ref{sec:concordance} is the paper's central limitation and should
govern how the score is used. $\Gamma$ is a loss-avoidance ranker with
essentially zero correlation to realized profit. On this instrument the two
are in tension by construction: the entropy factor buys partial-zone mass,
partial-zone mass is bought with credit, and credit is the dollars. Anyone
wishing to rank dollars should rank credit over width, accept a full-loss rate
above chance, and size accordingly; anyone wishing to rank survival should
rank $-\dent$ and accept thin credits. $\Gamma$ takes an interior position and
should be judged on whether that position is the one wanted, not on whether it
maximizes either axis.

It follows that the system has no dollar ranker. We regard finding one, or
proving that none exists within this instrument family, as the most valuable
open problem the configuration raises.

\subsection{Execution dominates everything}

The break-even fill of $0.92\times$ fitted fair value against measured fills
of $1.04\times$ is a $12$-point margin on a quantity estimated from $52$
trades in one account. Every headline number in Section~\ref{sec:results} is
conditional on it. Two structural facts make the margin more fragile than the
arithmetic suggests. The book trades near-the-money verticals at a $47\%$ win
rate and payoff odds near unity, so its profit is a $9\%$ residual of two
flows each twenty times its size; and selection is a within-day ranking on a
score whose own inputs move by five to seven percent between two chains of the
same session (Section~\ref{sec:reproducibility}). A systematic execution bias
of a few percent would not degrade the strategy, it would invert it.

The quote-reality gates of Section~\ref{sec:qualitygates} are the operational
response and should be read as a methodological point rather than an
implementation detail. Of $114{,}712$ candidates in a standard vendor frame,
$12.9\%$ had books that could not transact at all -- $14{,}768$ with a
non-positive immediate credit. A backtest of multi-leg option strategies that
does not test for this is pricing a substantial fraction of its book off
quotes that were never simultaneously available, and it will report an edge it
cannot capture.

\section{Conclusions}
\label{sec:conclusions}

The FCM configuration scores each candidate short-dated vertical credit spread
as a Kelly expected value under an empirically counted outcome measure,
discounted exponentially in the divergence of the \emph{market's} three-state
outcome measure from the discrete uniform. In the additive form of
Proposition~\ref{prop:additive} it ranks on
$g + k\,H(\Qbs)$: a growth criterion that pays for market uncertainty at $k$
nats of log-growth per nat of entropy.

Taking the entropy factor at the published uniform reference, with the
\emph{market's} measure as its argument rather than the trader's, makes it a
measurable property of the instrument. On this family it is an almost exact
monotone transform of the market's partial-zone probability, equivalently of
the market's expected move in strike widths, so the penalty asks a concrete
question about each candidate: is this a genuine three-state wager, or has the
market already priced the cushion between the strikes away? That question is
answerable from quoted prices alone, which is why the factor reproduces across
independent data sources at a rank correlation of $0.98$ where the estimated
win probability reproduces at $0.54$.

On $5{,}779$ in-sample and $814$ out-of-time trades the configuration returns
$+\$73{,}391$ and $+\$18{,}198$ from a \$10{,}000 bankroll at two contracts
per pick, with dollar Sharpe ratios of $1.26$ and $2.86$ and maximum drawdowns
of $-30.3\%$ and $-15.2\%$, booked at exactly the fitted fair credit and with
no assumed execution edge; a $146$-trade replay on live brokerage quotes is
positive on the same assumption. Ablation shows the composite score beating
each of its factors alone and beating random selection decisively, with its
own contribution concentrated in drawdown rather than return, and shows the
regime-gated bear sleeve -- eleven percent of the book -- carrying a
disproportionate share of its risk-adjusted performance.

The criterion's limitation is as clear as its mechanism. It ranks the
probability of losing and not the realized dollars, because on this instrument
the partial-zone cushion it rewards is bought with the credit that pays. That
tension is not a defect in the implementation; it is a property of selling
verticals, and naming it is the most useful thing this paper does for the work
that follows it.

\section*{Materials and Methods}

Backtests use vendor end-of-day option chains, 2020 through September 2026, on
S\&P~100 underlyings plus SPY, QQQ and IWM. Underlying close histories are
split-adjusted with per-event verification against the observed session
return. The realized-outcome measure $\Preal$ is counted on $252$ completed
DTE-matched windows per candidate against that candidate's own strikes, with
the window mean removed and a $\tfrac12$ pseudo-count per state, plus a
walk-forward own-gap drift refit each January. The market measure $\Qbs$ is
Black--Scholes $N(d_2)$ at the per-chain robust weighted quadratic implied
volatility fit. Earnings dates are the vetted union of a vendor calendar and
SEC Item 2.02 filings; ex-dividend dates and corporate actions come from a
public market-data endpoint. Live operation runs the same scoring code path on
brokerage snapshots at half-hour intervals, and the live replay of
Section~\ref{sec:live} rescores archived snapshots with the production
configuration. All computations are in Python with NumPy and pandas. Every
ablation in Section~\ref{sec:results} is evaluated from a single cached
scoring of the candidate pool, so all configurations share identical inputs,
identical settlements and identical fills.

\section*{Abbreviations}

\begin{tabular}{ll}
GEPO & generalized entropic portfolio optimization \\
REPO & return-entropy portfolio optimization \\
DEPO & discrete entropic portfolio optimization \\
GROUND & growth rate over uniform divergence \\
MVPO & mean-variance portfolio optimization \\
ATR & average true range \\
DTE & days to expiry \\
EV & expected value \\
IV & implied volatility \\
RV & realized volatility \\
SMA & simple moving average \\
\end{tabular}

\begin{thebibliography}{99}

\bibitem{markowitz1952} H.~Markowitz, ``Portfolio Selection,''
\emph{Journal of Finance}, 7(1):77--91, 1952.

\bibitem{shannon1948} C.~E.~Shannon, ``A Mathematical Theory of Communication,''
\emph{Bell System Technical Journal}, 27:379--423, 623--656, 1948.

\bibitem{kullback1951} S.~Kullback and R.~A.~Leibler, ``On Information and
Sufficiency,'' \emph{Annals of Mathematical Statistics}, 22(1):79--86, 1951.

\bibitem{kelly1956} J.~L.~Kelly, Jr., ``A New Interpretation of Information
Rate,'' \emph{Bell System Technical Journal}, 35(4):917--926, 1956.

\bibitem{sharpe1966} W.~F.~Sharpe, ``Mutual Fund Performance,''
\emph{Journal of Business}, 39(1):119--138, 1966.

\bibitem{thorp1971} E.~O.~Thorp, ``Portfolio Choice and the Kelly Criterion,''
\emph{Proceedings of the Business and Economics Section of the American
Statistical Association}, 215--224, 1971.

\bibitem{philippatos1972} G.~C.~Philippatos and C.~J.~Wilson, ``Entropy, Market
Risk, and the Selection of Efficient Portfolios,'' \emph{Applied Economics},
4(3):209--220, 1972.

\bibitem{blackscholes1973} F.~Black and M.~Scholes, ``The Pricing of Options and
Corporate Liabilities,'' \emph{Journal of Political Economy}, 81(3):637--654,
1973.

\bibitem{merton1973} R.~C.~Merton, ``Theory of Rational Option Pricing,''
\emph{Bell Journal of Economics and Management Science}, 4(1):141--183, 1973.

\bibitem{breeden1978} D.~T.~Breeden and R.~H.~Litzenberger, ``Prices of
State-Contingent Claims Implicit in Option Prices,'' \emph{Journal of Business},
51(4):621--651, 1978.

\bibitem{christensen1998} B.~J.~Christensen and N.~R.~Prabhala, ``The Relation
between Implied and Realized Volatility,'' \emph{Journal of Financial
Economics}, 50(2):125--150, 1998.

\bibitem{frittelli2000} M.~Frittelli, ``The Minimal Entropy Martingale Measure
and the Valuation Problem in Incomplete Markets,'' \emph{Mathematical Finance},
10(1):39--52, 2000.

\bibitem{follmer2002} H.~F\"ollmer and A.~Schied, ``Convex Measures of Risk and
Trading Constraints,'' \emph{Finance and Stochastics}, 6(4):429--447, 2002.

\bibitem{andersen2003} T.~G.~Andersen, T.~Bollerslev, F.~X.~Diebold, and
P.~Labys, ``Modeling and Forecasting Realized Volatility,''
\emph{Econometrica}, 71(2):579--625, 2003.

\bibitem{bakshikapadia2003} G.~Bakshi and N.~Kapadia, ``Delta-Hedged Gains and
the Negative Market Volatility Risk Premium,'' \emph{Review of Financial
Studies}, 16(2):527--566, 2003.

\bibitem{mackay2003} D.~J.~C.~MacKay, \emph{Information Theory, Inference and
Learning Algorithms}, Cambridge University Press, 2003.

\bibitem{coverthomas2006} T.~M.~Cover and J.~A.~Thomas, \emph{Elements of
Information Theory}, 2nd ed., Wiley-Interscience, 2006.

\bibitem{maccheroni2006} F.~Maccheroni, M.~Marinacci, and A.~Rustichini,
``Ambiguity Aversion, Robustness, and the Variational Representation of
Preferences,'' \emph{Econometrica}, 74(6):1447--1498, 2006.

\bibitem{hansensargent2008} L.~P.~Hansen and T.~J.~Sargent, \emph{Robustness},
Princeton University Press, 2008.

\bibitem{carrwu2009} P.~Carr and L.~Wu, ``Variance Risk Premiums,''
\emph{Review of Financial Studies}, 22(3):1311--1341, 2009.

\bibitem{bollerslev2009} T.~Bollerslev, G.~Tauchen, and H.~Zhou, ``Expected
Stock Returns and Variance Risk Premia,'' \emph{Review of Financial Studies},
22(11):4463--4492, 2009.

\bibitem{maclean2011} L.~C.~MacLean, E.~O.~Thorp, and W.~T.~Ziemba (eds.),
\emph{The Kelly Capital Growth Investment Criterion: Theory and Practice},
World Scientific, 2011.

\bibitem{zhou2013} R.~Zhou, R.~Cai, and G.~Tong, ``Applications of Entropy in
Finance: A Review,'' \emph{Entropy}, 15(11):4909--4931, 2013.

\bibitem{mercurio2020repo} P.~J.~Mercurio, Y.~Wu, and H.~Xie, ``An
Entropy-Based Approach to Portfolio Optimization,'' \emph{Entropy}, 22(3):332,
2020.

\bibitem{mercurio2020binary} P.~J.~Mercurio, Y.~Wu, and H.~Xie, ``Portfolio
Optimization for Binary Options Based on Relative Entropy,'' \emph{Entropy},
22(7):752, 2020.

\bibitem{mercurio2020gepo} P.~J.~Mercurio, Y.~Wu, and H.~Xie, ``Option
Portfolio Selection with Generalized Entropic Portfolio Optimization,''
\emph{Entropy}, 22(8):805, 2020.

\bibitem{muravyev2020} D.~Muravyev and N.~D.~Pearson, ``Options Trading Costs
Are Lower than You Think,'' \emph{Review of Financial Studies},
33(11):4973--5014, 2020.

\end{thebibliography}

\end{document}
